\section{Introduction}\label{sec:introduction}

The Coulomb interaction in two dimensions is logarithmic. An infinite
configuration of equal point charges against a uniform neutralizing
background has an infinite electrostatic energy, even in a bounded region
containing a charge. Subtracting the logarithmic self-energy at each charge
and then taking energy per unit area gives the renormalized energy studied
here. Its proposed minimizer is the triangular lattice.

\subsection{The energy and the main results}

For $R>1$, let $K_R=[-R,R]^2$. Fix smooth cutoffs
$\chi_R\in C_c^\infty(\R^2)$ satisfying
\begin{equation}\label{eq:cutoffs}
0\leq\chi_R\leq1,\qquad
\supp\chi_R\subset K_R,\qquad
\chi_R=1\text{ on }K_{R-1},\qquad
\sup_{R>1}\|\nabla\chi_R\|_\infty<\infty.
\end{equation}
Let $\Lambda\subset\R^2$ be locally finite and simple, so every point has
multiplicity one, and put $\nu_\Lambda=\sum_{p\in\Lambda}\delta_p$.
The admissible class $\mathcal A_1$ consists of locally integrable vector
fields $E$ such that
\begin{equation}\label{eq:admissible}
\operatorname{div}E=2\pi(\nu_\Lambda-1),\qquad
\curl E=0,\qquad
\sup_{R>1}\frac{\nu_\Lambda(B_R)}{|B_R|}<\infty,
\end{equation}
where the equations hold in distributions and $B_R$ is the disk of radius
$R$ centered at the origin. These conditions imply
$E\in L^q_{\mathrm{loc}}(\R^2;\R^2)$ for every $1\leq q<2$:
locally subtract the finitely many Coulomb singularities and a smooth
field with the prescribed constant divergence; the remaining field has
harmonic components. Section~\ref{sec:reduction} supplies the details.
Consequently the same class results if local $L^q$ integrability is required
for any one, or for all, of these exponents.

For $E\in\mathcal A_1$ define
\begin{align}
W(E,\chi_R)
&=\lim_{\eta\downarrow0}\left[
 \frac12\int_{\R^2\setminus\bigcup_{p\in\Lambda}B(p,\eta)}
       \chi_R|E|^2\dd x
 +\pi\log\eta\sum_{p\in\Lambda}\chi_R(p)\right],\label{eq:local-energy}\\
W(E)&=\limsup_{R\to\infty}\frac{W(E,\chi_R)}{|K_R|}.
\label{eq:whole-energy}
\end{align}
The local puncture limit exists and is finite. The outer limsup is
initially allowed to take either infinite value. The puncture limit is
taken first, and the cutoff layer has the fixed width in
\eqref{eq:cutoffs} throughout.

Identify $\R^2$ with $\mathbb C$ when convenient. The triangular lattice
of covolume one is
\begin{equation}\label{eq:triangular-lattice}
\Lambda_\triangle
=\sqrt{\frac{2}{\sqrt3}}
 \left(\Z(1,0)+\Z\left(\frac12,\frac{\sqrt3}{2}\right)\right).
\end{equation}
Let $h_\triangle$ be the mean-zero periodic distributional solution of
$-\Delta h_\triangle=2\pi(\delta_0-1)$ on
$\R^2/\Lambda_\triangle$, and set $E_\triangle=-\nabla h_\triangle$.

\begin{theorem}\label{thm:main}
For every cutoff family satisfying \eqref{eq:cutoffs} and every field
$E\in\mathcal A_1$,
\[
W(E)\geq W(E_\triangle).
\]
The triangular value is finite, and
$\inf_{E\in\mathcal A_1}W(E)=W(E_\triangle)$.
\end{theorem}

In particular $W(E)=-\infty$ is impossible in this class. The inequality
includes fields of energy $+\infty$. No separation, periodicity, or
Voronoi-cell bound is imposed on an arbitrary competitor.

The finite-volume statement used to prove Theorem~\ref{thm:main} has its
own exact normalization. For an integer $n\geq2$, write
$T_n=\R^2/(\sqrt n\,\Z^2)$. Given distinct points
$a_1,\ldots,a_n\in T_n$, let $h$ have mean zero and satisfy
$-\Delta h=2\pi(\sum_{i=1}^n\delta_{a_i}-1)$ on $T_n$. Define
\begin{equation}\label{eq:torus-energy}
\mathcal W_n(h)=\lim_{\eta\downarrow0}
\left[\frac12\int_{T_n\setminus\bigcup_{i=1}^n B(a_i,\eta)}
 |\nabla h|^2\dd x+\pi n\log\eta\right].
\end{equation}

\begin{theorem}\label{thm:torus}
For every integer $n\geq2$ and every configuration of $n$ distinct points
on $T_n$, the potential in \eqref{eq:torus-energy} satisfies
\[
\mathcal W_n(h)\geq nW(E_\triangle).
\]
\end{theorem}

This is a lower bound on square tori. It does not assert equality on any
particular square torus. Constant harmonic fields increase the energy by
$n|c|^2/2$; the required orthogonality follows by integrating the periodic
gradient, as proved in \eqref{red:harmonic}.

The planar minimum also determines the linear term in optimal
logarithmic energy for finitely many points on the sphere.

\begin{corollary}[Spherical logarithmic energy]\label{cor:bhs}
Let $S^2=\{y\in\R^3:\|y\|=1\}$ be the unit sphere. For each integer
$n\ge2$, define
\[
E_{\log}(n)=\min_{(y_1,\ldots,y_n)\in(S^2)^n}
 \left(-\sum_{\substack{1\le i,j\le n\\i\ne j}}
              \log\|y_i-y_j\|\right),
\]
where the sum is over ordered pairs, $\|\cdot\|$ is the Euclidean norm
in $\R^3$, $\log$ is the natural logarithm, and a collision has energy
$+\infty$. Then, as $n\to\infty$ through the integers,
\begin{equation}\label{eq:bhs-asymptotic}
E_{\log}(n)=\left(\frac12-\log2\right)n^2
             -\frac n2\log n+C_{\mathrm{BHS}}n+o(n),
\end{equation}
where
\begin{equation}\label{eq:bhs-constant}
C_{\mathrm{BHS}}
=2\log2+\frac12\log\!\left(\frac23\right)
       +3\log\!\left(\frac{\sqrt\pi}{\Gamma(1/3)}\right),
\end{equation}
and $\Gamma$ is Euler's gamma function.
\end{corollary}

This proves the $d=2$ case of the logarithmic-energy conjecture of
Brauchart, Hardin, and Saff~\cite[Conjecture~4]{BrauchartHardinSaff2012}.
B\'etermin and Sandier had already proved the existence of the linear
coefficient and shown that it equals \eqref{eq:bhs-constant} precisely
when the density-one triangular lattice minimizes their planar
renormalized energy~\cite[Theorem~1.5]{BeterminSandier2018}.
Section~\ref{sec:conclusion} checks the normalization and applies that
equivalence to Theorem~\ref{thm:main}. The conclusion is an asymptotic
value, not a construction of near-optimal point sets or an algorithm for
Smale's seventh problem.

\subsection{Context and prior work}

The finite-vortex renormalized energy of Bethuel, Brezis, and
H\'elein~\cite{BethuelBrezisHelein1994} is an antecedent of the
infinite-system energy considered here. Sandier and Serfaty introduced
the planar functional in their analysis of microscopic vortex patterns
in the Ginzburg--Landau model. They proved existence of a minimum,
approximation of its value by square-periodic fields, and triangular
optimality among Bravais lattices~\cite[Theorems~1 and~2]{SandierSerfaty2012}.
Their Conjecture~1 asks whether the triangular value is also the minimum
over the full admissible class. Theorem~\ref{thm:main} resolves this
conjecture positively. The current normalization in that paper has one
point per area $2\pi/m$; our density-one convention corresponds to its
background $m=2\pi$.

The lattice comparison belongs to a longer history. The minimization of
convergent Epstein zeta sums at fixed covolume was developed by
Rankin~\cite{Rankin1953} and Cassels~\cite{Cassels1959}, with an error
noted in~\cite{Cassels1963Corrigendum} and a repair supplied by
Ennola~\cite{Ennola1964Lemma} and simplified by Diananda~\cite{Diananda1964}.
Montgomery proved triangular minimality of the lattice theta function
for every positive Gaussian parameter~\cite[Theorem~1]{Montgomery1988}.
These are comparisons within the lattice class; the issue here is the
lower bound for arbitrary admissible configurations.

For the confining-potential class in their Theorem~2, Sandier and
Serfaty show that microscopic limits of finite Coulomb minimizers
concentrate on minimizers of the planar renormalized
energy~\cite[Theorem~2]{SandierSerfaty2015}.  Its minimum also determines
a next-order coefficient for optimal logarithmic energy on the
sphere~\cite[Theorem~1.5]{BeterminSandier2018}.  Petrache and Serfaty
prove a conditional implication from universal optimality to minimality
for Coulomb and Riesz renormalized energies in their configuration-energy
normalization~\cite[Theorem~2 and Corollary~4]{PetracheSerfaty2020}.
Their periodic comparison uses tori compatible with the reference
lattice, not arbitrary square tori, and their planar conclusion is
conditional on the Cohn--Kumar conjecture.  The companion universal-optimality
theorem supplies the Gaussian comparison used in that
implication~\cite[Theorem~1.1 and Corollary~1.3]{UniversalOptimality2026}.
The present proof treats
the logarithmic functional directly and uses no universal-optimality
result as an input.

Two geometric precedents help explain the finite-torus argument.
Screening a point charge by a uniform disk yields a separation argument
for planar confined logarithmic minimizers; see
Lieb, Rougerie, and Yngvason~\cite[Lemma~3 and Appendix~A]{LiebRougerieYngvason2018}.
Voronoi moment inequalities combine sharp cell comparisons with Euler's
bound on the average number of sides, as in Gruber's proof of a theorem
of Fejes T\'oth~\cite{Gruber1999}. Affine bounds on cell energies, sharp
at the regular hexagon, also appear in the Wasserstein model of
Bourne, Peletier, and Theil for cells with at least three sides and mass
$v\ge1.5\cdot10^{-4}$~\cite[Lemma~12]{BournePeletierTheil2014}.
Here the cell comparison must arise from a trial potential that is
continuous off the sites and has the prescribed logarithmic
singularities. Its nonradial correction and
the two scalar estimates of Section~\ref{sec:calibration} address that
requirement, including sectors whose perpendicular foot lies outside
the corresponding edge.

The whole-plane reduction below uses Sandier and Serfaty's local energy
estimates and periodic approximation. The finite-torus argument
provides the comparison with the triangular value. Throughout, the
conclusion concerns energy per unit area: compact changes to a
configuration can preserve this value, so the minimum statement alone
does not classify all minimizing configurations.


\subsection{How the comparison works}

The proof separates a finite-torus comparison from the passage to
arbitrary fields.  The finite comparison identifies the triangular
value.  Sandier and Serfaty's approximation theorem supplies periodic
fields approaching the unknown full-plane minimum; it does not identify
that minimum.

Section~\ref{sec:reduction} checks the conversion from that approximation
theorem to our normalization and fixed-width square cutoffs.  The main
point is to control the boundary error for an arbitrary field whose
energy limsup is bounded above, without assuming separation or density
one in advance. Local energy estimates give the needed charge density
and boundary-strip bound. Proposition~\ref{red:reduction} then transfers
any common lower bound on square tori to every field in $\mathcal A_1$,
including the possible infinite-energy cases.

We next prove the finite comparison.  On each square torus the point
energy attains a minimum.  A screened logarithmic potential and the
minimum principle separate the minimizing points, as shown in
Section~\ref{sec:geometry}.  Separation is used only for that minimum;
once it is bounded below, every other configuration is bounded below by
minimality.  This is why the main theorems impose no separation condition
on their competitors.

The lifted minimizer has convex Voronoi cells, the regions closer to
one site than to any other.  Each edge and its incident cell determine
a sector joining that edge to the site.  Counting incidences on the
torus gives at most six sectors per point on average.  The comparison is
not an estimate of the electrostatic energy of an isolated cell.
Instead, we construct a single trial potential $H$ on the torus. Its
sector formulas have matching traces and satisfy $H=-\log r+O(r^2)$
near each site. Thus $H$ and the Green potential $h$ have matching
logarithmic singularities, while the constant term is normalized to
zero in $H$. Proposition~\ref{dual:identity} gives the exact identity
\[
 \mathcal W_n(h)
 =\sum_S J(S)+\frac12\int_{T_n}|\nabla(h-H)|^2.
\]
Here $J(S)$ is the trial potential's sector contribution: its linear
background-interaction term minus one half of its renormalized Dirichlet
integral, as defined in \eqref{dual:sector-functional}. The identity
cancels the matching logarithmic singularities and retains the
nonnegative square exactly.

The construction of $H$ uses radial functions common to every sector.
At a shared Voronoi edge its trace depends only on the equal distances
to the two neighboring sites; along a ray to a vertex it depends only
on the vertex radius.  These observations glue the local formulas
continuously off the sites and give an $H^1$ function on every punctured
torus. The identity is proved for any compatible flat torus whose area
equals its number of sites, not just for square tori. A separate
one-site triangular torus supplies six regular sectors; there the chosen
potential has exactly the triangular Green gradient, so the square term
vanishes.

Section~\ref{sec:calibration} constructs the radial functions from the
triangular Green expansion and subtracts fixed multiples of sector area
and angle.  The desired bound has the form
\[
 J(S)\ge K+\lambda A(S)+\mu\theta(S),\qquad K<0.
\]
Areas and angles sum exactly to $n$ and $2\pi n$ on a square torus.
The incidence count then gives $n(6K+\lambda+2\pi\mu)$, and exact
calibration on the triangular torus identifies this with
$nW(E_\triangle)$.

The perpendicular from a site to a Voronoi edge need not lie on the
edge, so one endpoint angle of a sector can be negative.  A lower bound
for a scalar primitive alone does not control the resulting difference
of two primitive values.  A second estimate bounds the total negative
part of its derivative.  Together these estimates handle both ordinary
and signed sectors without imposing a regular-polygon hypothesis.

Section~\ref{sec:tails} proves summable analytic bounds for all omitted
lattice modes.  Section~\ref{sec:computation} first verifies the retained
finite core and then uses the analytic error bounds to enclose the full
functions and their derivatives.  Midpoint and interpolation errors
cover a compact rectangle between mesh points.  At the exact triangular
point, analytic stationarity and a positive Hessian bound establish the
non-strict minimum; an exact formula handles the unbounded exterior.
Section~\ref{sec:conclusion} combines these results with the geometry
and the whole-plane reduction, then derives Corollary~\ref{cor:bhs}.
Appendix~\ref{sec:scientific-support}
records two complementary methods: an independent coefficient recurrence
and sharper exact tail bounds.  They preserve additional quantitative
information without adding premises to the main comparison.
